\section{Actores, contexto y marco del proceso de requisitos}
\subsection{Actores principales y responsabilidades}
La identificaci\'on de interesados permite reconocer a quienes solicitan el proyecto, utilizan el producto o influyen en sus restricciones. Los papeles deben distinguirse: \emph{cliente} no es necesariamente \emph{usuario}, y el \emph{ingeniero de requisitos} no debe asumir decisiones comerciales que corresponden a la organizaci\'on \cite{pohl2025,swebok2024}. Para el PFC se propone entrevistar al representante del negocio, observar un flujo comercial autorizado y contrastar dudas con personal operativo. Todav\'ia no se registran entrevistas ejecutadas.

\begin{table}[!ht]
\centering\scriptsize
\caption{Cinco interesados iniciales del PFC; funciones propuestas para validar}
\label{tab:stakeholders}
\begin{tabularx}{\columnwidth}{@{}P{1.43cm}Y P{1.25cm}@{}}
\toprule Actor & Necesidad que conviene consultar & Rol en IR\\\midrule
Gerencia & Control de ventas, precios y alcance del proyecto. & Cliente y aprobador\\
Vendedor/a & B\'usqueda, cotizaci\'on y cierre de ventas. & Usuario primario\\
Encargado/a de inventario & Consistencia de existencias y devoluciones. & Usuario y experto de dominio\\
Comprador/a & Informaci\'on clara y sugerencias pertinentes. & Interesado externo\\
Equipo de tecnolog\'ia & Seguridad, operaci\'on y mantenibilidad. & Desarrollador / soporte\\
\bottomrule
\end{tabularx}
\end{table}

\subsection{Matriz de poder e inter\'es preliminar}
La clasificaci\'on de interesados se plantea como una \emph{herramienta de an\'alisis del equipo} inspirada en la necesidad de gestionar su participaci\'on en proyectos, tal como se aborda en PMBOK 7 \cite{pmbok2021}. No representa una evaluaci\'on autorizada por ELECTRORIOS y deber\'a revisarse con sus responsables.

\begin{table}[!ht]
\centering\scriptsize
\caption{Valoraciones provisionales de poder e inter\'es}
\label{tab:poder}
\begin{tabularx}{\columnwidth}{@{}P{1.35cm}ccY@{}}
\toprule Interesado & Poder & Inter\'es & Estrategia propuesta\\\midrule
Gerencia & Alto & Alto & Colaborar y solicitar aprobaci\'on\\
Vendedor & Medio & Alto & Demostraciones y revisi\'on de tareas\\
Inventario & Medio & Alto & Revisar reglas de stock\\
Comprador & Bajo & Alto & Contrastar claridad de informaci\'on\\
Soporte TIC & Medio & Medio & Consultar seguridad y factibilidad\\
\bottomrule
\end{tabularx}
\end{table}

\subsection{Delimitaci\'on del sistema y entidades externas}
El l\'imite preliminar comprende la interfaz de ventas, la administraci\'on del cat\'alogo, los registros de operaciones y el m\'odulo de recomendaci\'on. La distinci\'on entre entorno y sistema ayuda a no inventar integraciones ni responsabilidades \cite{pohl2025}. La Tabla~\ref{tab:interfaces} enumera \textbf{posibles} relaciones que el equipo estudiar\'a: el protocolo y la frecuencia no son datos observados.

\begin{table}[!ht]
\centering\scriptsize
\caption{Entidades externas e intercambios por confirmar}
\label{tab:interfaces}
\begin{tabularx}{\columnwidth}{@{}P{1.24cm}Y P{1.0cm}P{1.0cm}@{}}
\toprule Entidad & Datos previstos & Protocolo & Momento\\\midrule
Personal de ventas & Consulta, selecci\'on, confirmaci\'on & Web/app (propuesta) & Por venta\\
Cat\'alogo maestro & C\'odigo, nombre, precio, estado & Por definir & Por consulta\\
Sistema de inventario & Stock y movimientos & Por definir & Por venta\\
Servicio de pago & Confirmaci\'on e ID de transacci\'on, si aplica & Por definir & Por pago\\
Motor de IA externo & C\'odigos de productos y criterios, si se contrata & Por definir & Por sugerencia\\
\bottomrule
\end{tabularx}
\end{table}

\subsection{Roles del proceso y framework de IR}
El \textbf{analista de requisitos} organiza sesiones, registra decisiones y eval\'ua inconsistencias; el \textbf{cliente} valida prioridades comerciales; los \textbf{usuarios} describen tareas y comprueban prototipos; y los \textbf{desarrolladores} contrastan restricciones y viabilidad. Son papeles del proceso, que pueden recaer en varias personas, no nombres de puestos ya confirmados dentro del establecimiento \cite{iso29148,pohl2025}.

\begin{table}[!ht]
\centering\scriptsize
\caption{Aplicaci\'on propuesta del framework al PFC}
\label{tab:framework}
\begin{tabularx}{\columnwidth}{@{}P{1.13cm}Y P{1.4cm}@{}}
\toprule Actividad & Acci\'on planeada & Artefacto esperado\\\midrule
Elicitaci\'on & Preguntar y observar la b\'usqueda de productos. & Notas de campo\\
An\'alisis & Revisar conflictos de stock y permisos. & Matriz de conflictos\\
Especificaci\'on & Redactar requisitos identificados y trazables. & Lista de RF/RNF\\
Validaci\'on & Contrastar escenarios con el cliente. & Acta de revisi\'on\\
\midrule
Transversal & Versionar cambios y controlar fuentes. & Git, trazabilidad\\
Transversal & Medir claridad y viabilidad de requisitos. & Lista de control\\
\bottomrule
\end{tabularx}
\end{table}

\subsection{Primer mapa de necesidades y futuras verificaciones}
La tabla establece trazabilidad de \emph{propuestas}, no evidencia de aceptaci\'on. Una necesidad solo se considerar\'a validada tras comprobar su origen y la aprobaci\'on del responsable correspondiente \cite{iso29148}.
\begin{table}[!ht]
\centering\scriptsize
\caption{Trazabilidad preliminar que se debe validar}
\label{tab:traza}
\begin{tabularx}{\columnwidth}{@{}P{0.75cm}Y P{1.2cm}@{}}
\toprule Necesidad & Requisito a desarrollar & M\'etodo propuesto\\\midrule
N-01 & Consulta de producto, precio y stock. & Demostraci\'on\\
N-02 & Registro de venta con referencia \`unica. & Caso de prueba\\
N-03 & Recomendaci\'on justificable y descartable. & Revisi\'on usuario\\
N-04 & Restricciones de edici\'on seg\'un rol. & Prueba de permisos\\
\bottomrule
\end{tabularx}
\end{table}

\noindent\textit{Estado al cierre de este avance:} quedan pendientes el levantamiento autorizado, la verificaci\'on de protocolos, la lista completa de 15 requisitos, la matriz de evaluaci\'on individual y el mapa conceptual editable; estos elementos se desarrollar\'an en los siguientes d\'ias de trabajo.
